\chapter{INTRODUCTION}
\label{ch:introduction}

\section{Objectives}
\label{sec:intro-objectives}

The project set out to:
\begin{enumerate}[label=(\roman*)]
  \item Build an end-to-end testbed -- a benchmark FastAPI service, a
    deterministic, seeded open-loop traffic generator, Prometheus
    telemetry, and a Gymnasium-compatible environment wrapping a live
    Minikube cluster -- on which both an RL controller and HPA baselines
    could be run under identical conditions.
  \item Specify a frozen Markov Decision Process for the autoscaling
    problem, including explicit rules for what the agent should do when
    telemetry is missing, when there is genuinely no traffic, and a
    deterministic safety layer that the policy cannot override.
  \item Train a PPO agent live against this cluster for a pre-registered,
    full budget, and evaluate it against two HPA baselines (one using CPU
    and memory, one using CPU only, at the same target and replica bounds
    as every other arm) under a pre-registered protocol with a decision
    rule fixed before evaluation began.
  \item Report every result in full, including any run that ends in a
    safety termination rather than quietly excluding it, and diagnose any
    operational failure in enough depth to identify the specific, named
    design decision responsible.
  \item Report every design decision and every mid-course correction made
    along the way, including the ones that did not work, in an append-only
    log, so that the final results can be read alongside the full history
    of decisions that produced them rather than only the decisions that
    turned out to matter.
  \item Build a demonstration console that lets a reader watch the trained
    policy make decisions on deterministic, firewalled demo traffic,
    separate from every training and evaluation profile.
\end{enumerate}

\section{Background and Literature Survey}
\label{sec:intro-background}

\subsection{Background}

Cloud and container platforms let an operator change the amount of compute
capacity assigned to a service while that service is running. In a
Kubernetes cluster the usual unit of this change is the pod: a service is
deployed as a set of identical pod replicas behind a single logical
Deployment, and the number of replicas, not the size of any one pod, is the
main knob available for matching capacity to demand. Kubernetes itself grew
out of Google's internal Borg and Omega cluster managers, and the replica as
the unit of horizontal scaling is inherited directly from that lineage
\cite{burns2016borg}. Choosing how many replicas to run at any moment is a
genuine trade-off and not a solved problem: too few replicas and requests
queue behind an already busy pod, latency rises, and a service level
objective can be missed; too many replicas and CPU and memory sit idle while
still being paid for, whether that cost is measured in cloud billing or, as
in our case, in the finite CPU and memory of a single laptop running a
Minikube cluster. Surveys of cloud autoscaling methods going back more than a
decade treat this trade-off as the central question the field is trying to
answer, and group the many proposed answers into a small number of families:
fixed threshold rules, queueing-theoretic models, control-theoretic
feedback controllers, time-series forecasting, and, more recently,
reinforcement learning \cite{lorido2014review}\cite{qu2018survey}.

Kubernetes ships its own answer to this question in the form of the
Horizontal Pod Autoscaler, or HPA. The HPA periodically reads a utilization
metric for a Deployment's pods, usually average CPU utilization as a
percentage of the CPU each pod has requested, compares it to a target value
set by the operator, and resizes the Deployment proportionally so that, if
the controller's own simplifying assumptions held exactly, the new replica
count would bring utilization back to the target
\cite{k8s_hpa_docs}. This is a reactive rule: the HPA only responds after
utilization has already moved away from the target, it only has one number
to look at unless the operator configures several metrics explicitly, and it
has no notion of "this traffic pattern has happened before and here is what
happened last time" -- every decision cycle is evaluated from the current
reading alone. The HPA is nonetheless the default choice in almost every
real Kubernetes deployment we are aware of, for reasons this report takes
seriously rather than dismisses: it is simple to reason about, its behaviour
is predictable from its configuration, it has been hardened against years of
production use, and it fails in ways operators have learned to anticipate.
Any alternative autoscaler has to be judged against this baseline, not
against an idealized or hypothetical one.

Reinforcement learning (RL) offers a different kind of answer. Instead of a
fixed rule written by a human, an RL agent learns a policy, a mapping from
observed state to action, by interacting with an environment and receiving a
reward signal that scores how good each action turned out to be
\cite{sutton2018rl}. Applied to autoscaling, the hope is that an agent can
combine several signals at once (CPU, memory, request rate, latency),
notice trends in those signals rather than reacting only to their current
value, and weigh latency against resource cost through a single, explicit
reward function rather than through a hand-tuned threshold. Proximal Policy
Optimization (PPO) \cite{schulman2017ppo} is the specific RL algorithm we
use throughout this project; we describe it in plain terms in
Section~\ref{sec:intro-ppo-basics} below, and in full in
Chapter~\ref{ch:methodology}.

This report describes a system we built and ran end to end: a benchmark
service and deterministic traffic generator, a Gymnasium-compatible
environment wrapping a live Minikube cluster, an eleven-feature
multi-resource observation space, a seven-component reward, a deterministic
safety layer, a safety-aware tie-breaking rule, and a pre-registered
evaluation protocol comparing the trained PPO policy against two HPA
baselines. Every design decision in it is justified on its own technical
merits in Chapter~\ref{ch:methodology}, several of them against evidence
produced by the system's own training and evaluation runs -- for example,
the two sessions in which every ready replica of the deployment briefly fell
to zero (Section~\ref{sec:zero-ready-stops}) are part of this system's own
operating history, not a different design's failure, and they are exactly
why the deployment is never allowed to run with fewer than two replicas. The
system was trained and evaluated on a single-node Minikube cluster running
on a MacBook Air M1 with 8~GB of RAM, under the discipline of writing down,
before any evaluation run, what would count as a valid run and what would
count as one controller being "better" than another.

\subsection{Literature Survey}
\label{sec:related-work}

\subsubsection{Rule-Based and Predictive Autoscaling}

Threshold rules remain the dominant approach to autoscaling in practice
precisely because they are easy to configure and their behaviour is easy to
predict, a point made explicitly in the Lorido-Botran, Miguel-Alonso, and
Lozano survey of elastic cloud autoscaling techniques, which organizes prior
work into threshold-rule, queueing-model, control-theoretic, time-series
prediction, and reinforcement-learning families
\cite{lorido2014review}. Qu, Calheiros, and Buyya's survey of web
application autoscaling makes a related point about the difficulty of
choosing thresholds and cooldown periods that work well across varying
workloads, which is exactly the tuning problem our own HPA baselines had to
go through before they could be trusted as a fair comparison point (we
describe this tuning process, including the counterbalanced study that
selected the 260\% CPU target used by both our HPA baselines, in
Chapter~\ref{ch:methodology}) \cite{qu2018survey}. Taherizadeh and
Stankovski propose dynamic, multi-level rules for containerized applications
that adjust their own thresholds at runtime rather than using one fixed
value \cite{taherizadeh2019dynamic}, which is one step in the direction an
RL-based approach pushes further. At industrial scale, Google's Autopilot
system, described by Rzadca et al., sets both replica count and per-task
resource limits using statistics of past usage combined with a machine
learning recommender, and is deployed across a large fraction of Google's
internal clusters \cite{rzadca2020autopilot}; it is not evaluated against a
hand-tuned HPA-style baseline in the way this report evaluates PPO against
HPA, because at Google's scale there usually is no such baseline running
alongside it to compare against. Rattihalli, Govindaraju, Lu, and Tiwari
study vertical rather than horizontal autoscaling in Kubernetes, resizing a
pod's own resource requests instead of the replica count
\cite{rattihalli2019vertical}; we deliberately kept to the horizontal
setting throughout this project so that the HPA baseline would be a fair,
like-for-like comparison rather than a different kind of scaling decision
altogether.

\subsubsection{Learning-Based Autoscaling}

Mao, Alizadeh, Menache, and Kandula's early demonstration that a
policy-gradient agent can learn cluster job scheduling directly from
experience, without being told the scheduling heuristics a human expert
would use, is one of the papers most often cited as motivation for applying
RL to resource management problems generally \cite{mao2016resource}.
Horovitz and Arian apply tabular Q-learning with an explicit SLA objective
to cloud autoscaling \cite{horovitz2018qlearning}; Rossi, Nardelli, and
Cardellini apply model-based RL to combined horizontal and vertical scaling
of containers \cite{rossi2019horizontal}; Imdoukh, Ahmad, and Alfailakawi
predict load ahead of time with an LSTM network and scale in anticipation of
that prediction rather than reactively \cite{imdoukh2020lstm}; Toka,
Dobreff, Fodor, and Sonkoly combine several forecasters to scale Kubernetes
edge clusters \cite{toka2021edge}; and Qiu, Banerjee, Jha, Kalbarczyk, and
Iyer's FIRM system uses deep RL to reprovision resources for microservices
that are violating latency objectives \cite{qiu2020firm}. All of these
studies report results that are, to varying degrees, encouraging for
RL-based autoscaling; what distinguishes our project from essentially all of
them is where the training and evaluation data came from. Most of the work
above trains and/or evaluates in a simulator or against a recorded trace
replayed offline; our agent was trained with every single one of its
decision steps executed against a real Minikube cluster, at a real
30-second decision interval, with a real safety layer that could, and
twice did, stop a training session. We return to this contrast directly in
Table~\ref{tab:related-work-summary} and in the comparative analysis in
Section~\ref{sec:comparative-analysis}.

\subsubsection{Reliable Evaluation of Reinforcement Learning}

A third thread of related work is less about autoscaling specifically and
more about how to evaluate an RL method honestly at all. Henderson, Islam,
Bachman, Pineau, Precup, and Meger demonstrate that deep RL results can vary
substantially across random seeds and implementation details that papers
often do not report in enough detail to reproduce \cite{henderson2018matters}.
Agarwal, Schwarzer, Castro, Courville, and Bellemare argue for reporting
uncertainty explicitly when only a small number of runs are practically
possible, rather than reporting a single number as if it were exact
\cite{agarwal2021precipice} -- directly relevant to our own evaluation,
where cost and practicality limited us to three repeats per profile per arm
in the pre-registered campaign and exactly one repeat per cell in the
exploratory stress test, and where we report means and standard deviations
rather than single point estimates wherever more than one repeat exists.
Dulac-Arnold, Mankowitz, and Hester catalogue the obstacles that real
systems place in the way of RL that simulators do not: limited numbers of
samples, hard safety constraints, and partial or delayed observations
\cite{dulacarnold2019challenges}, all three of which appear directly in our
project (limited samples because every step costs thirty real seconds;
safety constraints enforced by our deterministic gating layer; and delayed,
sometimes missing, telemetry handled by the invalid-sample rules in
Section~\ref{sec:invalid-telemetry}). Garc\'ia and Fern\'andez survey safe
reinforcement learning broadly, including approaches that restrict the
action an agent is allowed to take using an external rule rather than
trusting the learned policy alone \cite{garcia2015safe}, which is the same
pattern our safety-aware tie-break rule and replica-bound gating follow.
Finally, Nosek, Ebersole, DeHaven, and Mellor's argument for pre-registration
as a defence against outcome-driven analysis choices, originally made for
psychology experiments, is the methodological practice we adopted most
directly: our evaluation workload, our metrics, our retry rules, and our
decision rule for declaring a controller "better" were all fixed and hashed
before any evaluation episode ran \cite{nosek2018preregistration}.

\subsubsection{Summary Comparison of Closest Prior Work}

Table~\ref{tab:related-work-summary} summarizes how the studies above
compare along the dimensions most relevant to this project: whether training
and evaluation happened on a live system or in simulation/on a trace,
which resource signal(s) the method used, and whether the baseline
comparison, if any, was specified before results were known.

\begin{table}[H]
\centering
\caption[Comparison with closest prior work]{Summary comparison of the closest prior work surveyed in this chapter.}
\label{tab:related-work-summary}
\small
\begin{tabular}{p{2.6cm}p{2.6cm}p{2.3cm}p{2.3cm}p{3.5cm}}
\toprule
Study & Method family & Live or simulated/trace & Signal(s) used & How this project differs \\
\midrule
Mao et al.\ \cite{mao2016resource} & Policy-gradient RL & Simulated & Job/cluster scheduling features & We train every step live against Kubernetes, not a scheduling simulator. \\
Horovitz \& Arian \cite{horovitz2018qlearning} & Tabular Q-learning & Not specified as live & SLA-derived state & We use a continuous 11-feature observation and PPO, not a tabular Q-table. \\
Rossi et al.\ \cite{rossi2019horizontal} & Model-based RL & Simulated/model & CPU, container count & We scale horizontally only and train model-free, live. \\
Imdoukh et al.\ \cite{imdoukh2020lstm} & LSTM forecasting + rule & Trace-driven & Predicted load & We do not forecast; the agent reacts to current and trend features. \\
Toka et al.\ \cite{toka2021edge} & Forecast ensemble & Trace-driven & Multiple forecasters & Our comparison baseline is a tuned HPA, not an ensemble of predictors. \\
Qiu et al.\ \cite{qiu2020firm} & Deep RL & Production-scale, not a single baseline & Microservice latency SLOs & We run a pre-registered, held-out, repeated-measures comparison against HPA; FIRM does not report this kind of baseline comparison. \\
This project & On-policy PPO & Live, single-node Minikube, no simulator & 11 signals: CPU, memory, request rate, latency, replica state, all with short-horizon trends & Pre-registered evaluation protocol and decision rule fixed before any evaluation run; every result, successful or not, reported in full. \\
\bottomrule
\end{tabular}
\end{table}

\section{Problem Statement}
\label{sec:problem-statement}

Can a PPO agent, trained in real time against a live, single-node Kubernetes
cluster rather than against a simulator or a replayed trace, match or
improve on a correctly tuned Horizontal Pod Autoscaler when both are
evaluated on traffic that neither has seen during training, under a
decision rule that is written down before any evaluation result exists? The
honest answer this report gives is that the evidence does not support
declaring either kind of controller the winner: at the light evaluation
loads our pre-registered protocol exercised, the primary latency objective
does not discriminate between controllers at all, and on the cost metric
that does discriminate, no single arm is cheapest across every workload
profile we tested.

\section{Motivation}

Three considerations motivated taking on this question rather than, for
example, building a simulator-based RL autoscaler and reporting simulated
results, which is the much more common approach in the literature we survey
in Section~\ref{sec:related-work}. First, a simulator necessarily encodes
assumptions about how a real cluster behaves under load -- how quickly a new
pod becomes ready, how CPU throttling interacts with request latency, how
Prometheus scrape timing interacts with a 30-second decision interval -- and
those assumptions are exactly the things we did not trust ourselves to get
right without first seeing them happen on a real cluster. Training live
forces every one of those assumptions to be confronted directly, at the
cost of much slower data collection: our training campaign took 12,288 real
decision steps, each one tied to a real 30-second interval, which is a
meaningfully different experience from training an agent for the same
number of steps in a simulator that can run far faster than real time.
Second, much of the RL-for-autoscaling literature we found sets up its
baseline comparison after the RL method's results are already known, which
creates an obvious risk that the baseline is less carefully tuned than it
should be, or that the evaluation workload is chosen because it happens to
favour the proposed method. We wanted to avoid that risk by writing down, in
advance, exactly which traffic profiles would be used for evaluation, which
metrics would be measured, and what would count as one controller being
better than another, before running a single evaluation episode. This
discipline, borrowed from the pre-registration movement in experimental
psychology \cite{nosek2018preregistration}, shapes the whole structure of
Chapter~\ref{ch:methodology} and Chapter~\ref{ch:results}. Third, and more
simply, we wanted to know, for our own understanding of the problem, what
actually goes wrong when an RL autoscaler is trained and evaluated live
rather than in simulation, including the operational failures (a safety
termination from zero ready replicas, a training session that had to be
retried, an evaluation run blocked by a leftover HPA object) that a purely
simulated project would never encounter and that a report written only
about the successful runs would hide.

\section{Basic Concepts}
\label{sec:intro-basics}

This section introduces, briefly and in plain language, the concepts that
the rest of the report assumes. Readers already familiar with Kubernetes and
reinforcement learning can skip to Section~\ref{sec:existing-system}.

\subsection{Pods, Replicas, and the HPA Algorithm}

A \textbf{pod} is the smallest deployable unit in Kubernetes: one or more
containers that are scheduled together onto the same node and share a
network namespace. A \textbf{Deployment} manages a set of identical pods,
called \textbf{replicas}, and is responsible for keeping the number of
healthy, \textbf{ready} replicas equal to a desired count, restarting or
replacing pods that fail their health checks. \textbf{Horizontal scaling}
means changing this desired replica count; it is distinct from
\textbf{vertical scaling}, which changes how much CPU or memory each
individual pod is allowed to use. This report is about horizontal scaling
only, and every controller we compare -- the PPO agent and both HPA
configurations -- operates by changing the same Deployment's replica count,
never by resizing a pod.

The HPA controller runs a control loop, by default every fifteen seconds,
that: (1) reads the current value of its configured metric (for example,
average CPU utilization across all ready pods in the target Deployment) from
the Metrics Server or, when configured for custom/external metrics, from
Prometheus via the Prometheus Adapter; (2) computes a desired replica count
as $\mathtt{ceil}(\text{currentReplicas} \times \text{currentMetricValue} /
\text{targetMetricValue})$, with safeguards against scaling on a stale or
missing metric; (3) applies configured stabilization windows and scale
velocity policies that slow down how quickly it is allowed to scale up or
down, to avoid oscillation; and (4) patches the Deployment's replica count
if the computed value differs from the current one and the safeguards allow
it \cite{k8s_hpa_docs}. When an HPA is configured against more than one
metric, as our HPA CPU+Memory arm is, it computes a desired replica count
from each metric independently and takes the largest of the two, so that
the metric under more pressure always wins.

\subsection{Prometheus and Metrics Server}

Kubernetes ships a lightweight \textbf{Metrics Server} that samples CPU and
memory usage per pod on a short interval (we verified a 15-second resolution
live in our cluster before every HPA run) and exposes it through the
Kubernetes Metrics API; this is what the stock HPA reads for CPU/memory
percentage targets. \textbf{Prometheus} is a separate, general-purpose
time-series monitoring system that scrapes metrics endpoints exposed by
applications (our benchmark service exposes one at \texttt{/\allowbreak \allowbreak metrics}) and
stores them so that they can be queried later with PromQL, Prometheus's own
query language. Our RL agent's observation builder and our evaluation
analysis scripts both read from Prometheus rather than from the Metrics
Server, because Prometheus is the only one of the two that exposes
application-level signals such as request rate, request latency
percentiles, and error rate, none of which the Metrics Server knows about.
We installed Prometheus (and, only while its dashboards were actively being
used, Grafana) via the \texttt{kube-prometheus-stack} Helm chart; Grafana
itself played no role in any controller's decision loop and we scaled it to
zero replicas during both the training campaign and the evaluation campaign
so that it could not compete with the benchmark application for the
cluster's limited CPU and memory.

\subsection{Reinforcement Learning and the Markov Decision Process}

A reinforcement learning problem is usually formalized as a Markov Decision
Process (MDP): a set of states $s$, a set of actions $a$, a transition
behaviour (here, the real Kubernetes cluster and the traffic it is serving,
not a model of them), and a reward function $r(s,a)$ that scores each
transition. An \textbf{agent} chooses actions according to a
\textbf{policy} $\pi(a \mid s)$ and its goal is to choose a policy that
maximizes the expected sum of (discounted) future rewards, not just the
immediate reward of the next action. The \textbf{discount factor} $\gamma
\in [0,1)$ controls how much the agent cares about reward further in the
future relative to reward right now; we use $\gamma = 0.99$ throughout,
meaning the agent is encouraged to think a reasonable distance ahead rather
than acting greedily on the very next step. In our case, one
\textbf{decision step} corresponds to one 30-second control interval: the
agent observes the cluster's current state, chooses to remove one replica,
hold, or add one replica, and thirty seconds later observes the result and
receives a reward before choosing again. An \textbf{episode} is one
complete run of a traffic profile from start to end, after which the
environment resets.

\subsection{Proximal Policy Optimization in Plain Terms}
\label{sec:intro-ppo-basics}

PPO is an on-policy, actor-critic RL algorithm: it learns both a policy (the
\textit{actor}, which chooses actions) and a value function (the
\textit{critic}, which estimates how good a given state is expected to be),
and it updates both using data collected by the current policy itself,
discarding that data once the policy has moved on \cite{schulman2017ppo}.
The central idea that gives PPO its name is a \textit{clipped} objective:
when PPO updates the policy from a batch of collected experience, it limits
how far any single update is allowed to push the policy's action
probabilities away from what they were when the experience was collected.
This matters in our setting more than it might in a simulated one, because
every single step of experience in this project costs thirty real seconds
and a live interaction with the cluster -- there is no option to collect a
huge batch of experience cheaply and then take a large, aggressive update
step, the way one might in a fast simulator. PPO's clipping keeps each
update's step size small and stable, which is one reason it is a common
default choice for problems where data is expensive, slow, or must be
collected from a real system. We used the Stable-Baselines3 implementation
of PPO \cite{raffin2021sb3} throughout, built on PyTorch, wrapped in a
Gymnasium-compatible environment \cite{towers2024gymnasium} that we wrote to
represent our live Kubernetes cluster as an RL environment. The full set of
PPO hyperparameters we used is given in Table~\ref{tab:ppo-hyperparams} in
Chapter~\ref{ch:methodology}.

\section{Existing System and Its Limitations}
\label{sec:existing-system}

The existing system, in the sense of "what a Kubernetes operator would use
today without this project," is the stock Horizontal Pod Autoscaler
configured against one or more Metrics-Server- or Prometheus-adapter-backed
metrics. Its strengths are real: it is simple, its behaviour follows
directly from its configuration, and it is battle-tested. Its limitations,
which this project's motivation rests on, are equally real. It is purely
reactive: it cannot act on a trend, only on a current value, so a traffic
ramp that is clearly accelerating gets the same treatment as one that has
just stabilized at the same instantaneous value. It generally optimizes a
single metric (or the maximum over a small, independently-evaluated set of
metrics) rather than a combined objective that trades latency against
resource cost explicitly. Its thresholds and stabilization windows must be
hand-tuned for a given workload and, once tuned, do not adapt if the
workload's character changes. And it has no way to use what happened during
a previous, similar traffic episode to inform the current decision -- every
decision cycle starts from nothing but the current metric reading.

\section{Proposed Approach}

We propose, and report on, replacing the HPA's decision rule with a learned
policy while keeping everything else about the deployment -- the benchmark
service, its resource requests and limits, its health probes, the Minikube
cluster itself -- identical across every controller being compared. The
learned policy observes eleven normalized signals derived from Prometheus
telemetry (Chapter~\ref{ch:methodology}, Table~\ref{tab:observation-features}),
is trained with PPO against the live cluster, and its raw action is passed
through a deterministic safety layer before it ever reaches the Kubernetes
API, so that the learned part of the system can never, by itself, violate a
replica bound, bypass a cooldown, or act on telemetry that has been flagged
as invalid. We evaluate this policy against two HPA configurations --
one using CPU and memory, one using CPU only -- both at the same 260\% CPU
target and 2-8 replica bounds, under a protocol, decision rule, and set of
held-out traffic profiles that were all fixed and recorded before any
evaluation episode was run.

\section{Organization of the Report}

Chapter~\ref{ch:methodology} describes the proposed system and the full
methodology: the testbed, the benchmark service, the traffic generator, the
MDP formulation (observation, action space, safety layer, reward), the
standards the system relies on, the software and hardware used, the
evaluation protocol, and the demonstration interface. Chapter~\ref{ch:cost}
gives a cost analysis of the hardware and software used. Chapter~\ref{ch:results}
reports the training campaign, the held-out evaluation, the exploratory
stress test, the demonstration interface, a comparative analysis against
the related work surveyed in this chapter, and the project's limitations.
Chapter~\ref{ch:conclusion} concludes and discusses future work.
Chapter~\ref{ch:appendix} is an appendix containing verbatim code excerpts,
configuration files, a summary of the design-correction log, and
supplementary data tables. Chapter~\ref{ch:references} lists references in
IEEE numbered style, followed by the biodata of each student author.
